% GENERATED FILE. Edit canonical lessons, then run npm run generate:books.
% canonical-source-hash: fnv1a64:4313edf48c9782c5
% canonical-lessons: MR-C439-cakrivadal, MR-C439-vavtal, MR-C439-gadgadat, MR-C439-ghongavne, MR-C439-kadadne

\chapter{Storms and Thunder}
\label{ch:mr-storms-and-thunder}
\hlchaptermodality{439}

\begin{chapteropening}
\textbf{By the end of this chapter:} \emph{I can talk about cyclones, whirlwinds and thunder, and say the wind howled and lightning cracked.}

Complete the last lesson of chapter 439: I can talk about cyclones, whirlwinds and thunder, and say the wind howled and lightning cracked.
\end{chapteropening}

\section[\texorpdfstring{\mr{चक्रीवादळ} (cakrīvādaḷ)}{चक्रीवादळ (cakrīvādaḷ)}]{\mr{चक्रीवादळ} (cakrīvādaḷ) — a cyclone}

\label{lesson:MR-C439-cakrivadal}



\begin{quote}
Before the new one: say the Marathi for to pour down, then the Marathi for to get wet.
\end{quote}

\subsection*{You'll want to know: \mr{चक्रीवादळ}}
\textbf{\mr{चक्रीवादळ}} — \emph{cakrīvādaḷ} — \textquotedblleft{}a cyclone\textquotedblright{}.

\begin{quote}
\mr{चक्रीवादळामुळे} \mr{गावात} \mr{खूप} \mr{पडझड} \mr{झाली}.
\end{quote}

\begin{quote}
\emph{cakrīvādaḷāmuḷe gāvāt khūp paḍjhaḍ jhālī.}
\end{quote}

\begin{quote}
The cyclone brought down a lot in the village.
\end{quote}

Neuter.

\mr{चक्री}, whirling, and \mr{वादळ}, a storm.

\begin{cousinweb}
\textbf{\mr{चक्रीवादळ}} (\emph{cakrīvādaḷ}) is built from \emph{cakrī}, \textquotedblleft{}whirling\textquotedblright{}, on Sanskrit \emph{cakra}, \textquotedblleft{}a wheel\textquotedblright{} (CDIAL 4538), and \textbf{\mr{वादळ}}, \textquotedblleft{}a storm\textquotedblright{}.
\end{cousinweb}

\begin{grammarlens}[title={What you've built}]
One more word for getting around.
\end{grammarlens}

\subsection*{Your turn}
\begin{itemize}
\raggedright
  \item \emph{Say it:} \emph{cakrīvādaḷ}
  \item \emph{Say it:} \emph{cakrīvādaḷ}, once more
  \item \emph{Say it:} \emph{cakrīvādaḷāmuḷe gāvāt khūp paḍjhaḍ jhālī.} — The cyclone brought down a lot in the village.
  \item \emph{Recall:} say the Marathi for to pour down, then the Marathi for to get wet, then say \emph{cakrīvādaḷ} again
\end{itemize}

\begin{tcolorbox}[breakable,colback=teal!4,colframe=teal!35!black,title={Before you move on}]
What does \mr{चक्रीवादळ} mean? (\textquotedblleft{}A cyclone\textquotedblright{}.) Say it once more.
\end{tcolorbox}
\section[\texorpdfstring{\mr{वावटळ} (vāvṭaḷ)}{वावटळ (vāvṭaḷ)}]{\mr{वावटळ} (vāvṭaḷ) — a whirlwind}

\label{lesson:MR-C439-vavtal}



\begin{quote}
Before the new one: say the Marathi for to get wet, then the Marathi for a cyclone.
\end{quote}

\subsection*{You'll want to know: \mr{वावटळ}}
\textbf{\mr{वावटळ}} — \emph{vāvṭaḷ} — \textquotedblleft{}a whirlwind\textquotedblright{}.

\begin{quote}
\mr{शेतात} \mr{अचानक} \mr{वावटळ} \mr{उठली}.
\end{quote}

\begin{quote}
\emph{śetāt acānak vāvṭaḷ uṭhlī.}
\end{quote}

\begin{quote}
Suddenly a whirlwind rose in the field.
\end{quote}

Feminine.

\mr{वावटळ} \mr{उठणे}, a whirlwind rises.

\begin{cousinweb}
\textbf{\mr{वावटळ}} (\emph{vāvṭaḷ}) is a Marathi word for a spinning pillar of wind and dust; its origin is not settled.
\end{cousinweb}

\begin{grammarlens}[title={What you've built}]
One more word for getting around.
\end{grammarlens}

\subsection*{Your turn}
\begin{itemize}
\raggedright
  \item \emph{Say it:} \emph{vāvṭaḷ}
  \item \emph{Say it:} \emph{vāvṭaḷ}, once more
  \item \emph{Say it:} \emph{śetāt acānak vāvṭaḷ uṭhlī.} — Suddenly a whirlwind rose in the field.
  \item \emph{Recall:} say the Marathi for to get wet, then the Marathi for a cyclone, then say \emph{vāvṭaḷ} again
\end{itemize}

\begin{tcolorbox}[breakable,colback=teal!4,colframe=teal!35!black,title={Before you move on}]
What does \mr{वावटळ} mean? (\textquotedblleft{}A whirlwind\textquotedblright{}.) Say it once more.
\end{tcolorbox}
\section[\texorpdfstring{\mr{गडगडाट} (gaḍgaḍāṭ)}{गडगडाट (gaḍgaḍāṭ)}]{\mr{गडगडाट} (gaḍgaḍāṭ) — a rumble of thunder}

\label{lesson:MR-C439-gadgadat}



\begin{quote}
Before the new one: say the Marathi for a cyclone, then the Marathi for a whirlwind.
\end{quote}

\subsection*{You'll want to know: \mr{गडगडाट}}
\textbf{\mr{गडगडाट}} — \emph{gaḍgaḍāṭ} — \textquotedblleft{}a rumble of thunder\textquotedblright{}.

\begin{quote}
\mr{रात्री} \mr{ढगांचा} \mr{गडगडाट} \mr{ऐकू} \mr{आला}.
\end{quote}

\begin{quote}
\emph{rātrī ḍhagāñcā gaḍgaḍāṭ aikū ālā.}
\end{quote}

\begin{quote}
At night the rumble of thunder could be heard.
\end{quote}

Masculine.

\mr{ऐकू} \mr{येणे}, to be heard: the sound is the subject.

\begin{cousinweb}
\textbf{\mr{गडगडाट}} (\emph{gaḍgaḍāṭ}) is an echo word that imitates rolling thunder; its origin is not settled.
\end{cousinweb}

\begin{grammarlens}[title={What you've built}]
One more word for getting around.
\end{grammarlens}

\subsection*{Your turn}
\begin{itemize}
\raggedright
  \item \emph{Say it:} \emph{gaḍgaḍāṭ}
  \item \emph{Say it:} \emph{gaḍgaḍāṭ}, once more
  \item \emph{Say it:} \emph{rātrī ḍhagāñcā gaḍgaḍāṭ aikū ālā.} — At night the rumble of thunder could be heard.
  \item \emph{Recall:} say the Marathi for a cyclone, then the Marathi for a whirlwind, then say \emph{gaḍgaḍāṭ} again
\end{itemize}

\begin{tcolorbox}[breakable,colback=teal!4,colframe=teal!35!black,title={Before you move on}]
What does \mr{गडगडाट} mean? (\textquotedblleft{}A rumble of thunder\textquotedblright{}.) Say it once more.
\end{tcolorbox}
\section[\texorpdfstring{\mr{घोंगावणे} (ghoṅgāvṇe)}{घोंगावणे (ghoṅgāvṇe)}]{\mr{घोंगावणे} (ghoṅgāvṇe) — to howl; to roar}

\label{lesson:MR-C439-ghongavne}



\begin{quote}
Before the new one: say the Marathi for a whirlwind, then the Marathi for a rumble of thunder.
\end{quote}

\subsection*{You'll want to know: \mr{घोंगावणे}}
\textbf{\mr{घोंगावणे}} — \emph{ghoṅgāvṇe} — \textquotedblleft{}to howl; to roar\textquotedblright{}.

\begin{quote}
\mr{रात्रभर} \mr{वारा} \mr{घोंगावत} \mr{होता}.
\end{quote}

\begin{quote}
\emph{rātrabhar vārā ghoṅgāvat hotā.}
\end{quote}

\begin{quote}
The wind howled all night.
\end{quote}

Past (he …): \mr{तो} \mr{घोंगावला}, \emph{to ghoṅgāvlā}.

For strong wind and for swarming bees.

\begin{cousinweb}
\textbf{\mr{घोंगावणे}} (\emph{ghoṅgāvṇe}) is an echo verb for the roar of wind or the drone of bees; its origin is not settled.
\end{cousinweb}

\begin{grammarlens}[title={What you've built}]
One more word for getting around.
\end{grammarlens}

\subsection*{Your turn}
\begin{itemize}
\raggedright
  \item \emph{Say it:} \emph{ghoṅgāvṇe}
  \item \emph{Say it:} \emph{ghoṅgāvṇe}, once more
  \item \emph{Say it:} \emph{rātrabhar vārā ghoṅgāvat hotā.} — The wind howled all night.
  \item \emph{Recall:} say the Marathi for a whirlwind, then the Marathi for a rumble of thunder, then say \emph{ghoṅgāvṇe} again
\end{itemize}

\begin{tcolorbox}[breakable,colback=teal!4,colframe=teal!35!black,title={Before you move on}]
What does \mr{घोंगावणे} mean? (\textquotedblleft{}To howl; to roar\textquotedblright{}.) Say it once more.
\end{tcolorbox}
\section[\texorpdfstring{\mr{कडाडणे} (kaḍāḍṇe)}{कडाडणे (kaḍāḍṇe)}]{\mr{कडाडणे} (kaḍāḍṇe) — to crackle; to crash}

\label{lesson:MR-C439-kadadne}



\begin{quote}
Before the new one: say the Marathi for a rumble of thunder, then the Marathi for to howl.
\end{quote}

\subsection*{You'll want to know: \mr{कडाडणे}}
\textbf{\mr{कडाडणे}} — \emph{kaḍāḍṇe} — \textquotedblleft{}to crackle; to crash\textquotedblright{}.

\begin{quote}
\mr{आकाशात} \mr{वीज} \mr{कडाडली}.
\end{quote}

\begin{quote}
\emph{ākāśāt vīj kaḍāḍlī.}
\end{quote}

\begin{quote}
Lightning crashed across the sky.
\end{quote}

Past (he …): \mr{तो} \mr{कडाडला}, \emph{to kaḍāḍlā}.

Lightning \mr{कडाडते}; a fire can \mr{कडाडणे} too.

\begin{cousinweb}
\textbf{\mr{कडाडणे}} (\emph{kaḍāḍṇe}) is inherited, from the family of Sanskrit \emph{kaṭakaṭā}, \textquotedblleft{}a grating sound\textquotedblright{}, where CDIAL lists Marathi \emph{kaḍāḍ}, \textquotedblleft{}crackling\textquotedblright{} (CDIAL 2632).
\end{cousinweb}

\begin{grammarlens}[title={What you've built}]
One more word for getting around.
\end{grammarlens}

\subsection*{Your turn}
\begin{itemize}
\raggedright
  \item \emph{Say it:} \emph{kaḍāḍṇe}
  \item \emph{Say it:} \emph{kaḍāḍṇe}, once more
  \item \emph{Say it:} \emph{ākāśāt vīj kaḍāḍlī.} — Lightning crashed across the sky.
  \item \emph{Recall:} say the Marathi for a rumble of thunder, then the Marathi for to howl, then say \emph{kaḍāḍṇe} again
\end{itemize}

\begin{tcolorbox}[breakable,colback=teal!4,colframe=teal!35!black,title={Before you move on}]
What does \mr{कडाडणे} mean? (\textquotedblleft{}To crackle; to crash\textquotedblright{}.) Say it once more.
\end{tcolorbox}
